\documentclass[11pt,a4paper]{article}
\usepackage[margin=1in]{geometry}
\usepackage{times}
\usepackage{latexsym}
\usepackage{graphicx}
\usepackage{amsmath}
\usepackage{amssymb}
\usepackage{booktabs}
\usepackage{multirow}
\usepackage{xcolor}
\usepackage[hyphens]{url}
\usepackage{tabularx}
\usepackage{subcaption}
\usepackage{natbib}
\usepackage[colorlinks=true,citecolor=blue,urlcolor=blue]{hyperref}

\newcommand{\SJ}{\textsc{SJ}}
\newcommand{\FK}{\textsc{FK}}
\newcommand{\LP}{\textsc{LogProb}}
\newcommand{\EM}{\textsc{EM}}
\newcommand{\EV}{\textsc{EntVote}}

\title{A Systematic Comparison of Confidence Signals for Test-Time Scaling\\in Structured Information Extraction}

\author{Anonymous}

\begin{document}
\maketitle

% ==============================================================================
\begin{abstract}
Test-time scaling (TTS) through sampling and selection offers a promising approach to improving LLM outputs, but requires reliable confidence signals.
For structured prediction tasks such as information extraction (IE), it remains unclear which confidence signals---model-internal (sequence log-probability), surface agreement (Fleiss' Kappa, exact-match voting), or structural similarity (soft Jaccard)---best correlate with output quality.
We present the first systematic comparison of four confidence signal families for IE test-time scaling, evaluated across three NER datasets (SciERC, CoNLL-2003, WNUT-17), one RE dataset (SciERC), two model families (Qwen3-8B, LLaMA-3.1-8B), and multiple evaluation metrics ($\rho$, $\tau$, AUROC, ECE).
We find three key results:
(1)~structural consistency (soft Jaccard) consistently outperforms surface consistency (Fleiss' Kappa, exact-match) on NER ($\rho$=0.30--0.41 vs.\ 0.27--0.32), establishing a hierarchy among consistency-based signals;
(2)~sequence log-probability is a surprisingly strong zero-cost baseline ($\rho$=0.30--0.36), competitive with soft Jaccard on well-calibrated tasks but capturing fundamentally different information;
(3)~all confidence signals degrade on relation extraction (full-set $\rho < 0.28$), revealing a task-complexity boundary governed by compound error propagation---yet soft Jaccard recovers strong signal ($\rho$=0.37) within the model's capability range.
Best-of-N selection yields +2.3--2.8pp F1 over greedy decoding across NER datasets, demonstrating the practical value of TTS for structured prediction.
\end{abstract}

% ==============================================================================
\section{Introduction}

Test-time scaling (TTS)---the strategy of generating multiple candidate outputs and selecting the best one---has emerged as a practical approach to improving LLM performance without retraining \cite{wang2023selfconsistency,brown2024large}.
The key challenge is the \emph{selection problem}: identifying which candidate is most likely correct requires a reliable confidence signal.

For reasoning tasks, self-consistency via majority voting provides a natural solution: the most frequent answer among sampled outputs is selected \cite{wang2023selfconsistency}.
However, structured prediction tasks such as information extraction (IE) pose unique challenges.
IE outputs---entity sets, relation triplets---have internal structure that makes exact-match voting impractical: even high-quality outputs rarely match exactly due to minor span boundary variations.

Several alternatives exist: model-internal confidence (sequence log-probability), surface-level agreement metrics (Fleiss' Kappa), and structural similarity measures (soft Jaccard over entity/relation sets).
Yet no systematic comparison of these confidence signal families exists for IE test-time scaling.
This gap leaves practitioners without principled guidance on which signal to use, and whether the additional inference cost of multi-sample consistency metrics is justified over zero-cost model confidence.

We address this gap with a comprehensive empirical study comparing four families of confidence signals for IE TTS:
\begin{itemize}
    \item \textbf{Model-internal}: sequence log-probability (zero additional cost)
    \item \textbf{Surface agreement}: Fleiss' Kappa and exact-match voting rate
    \item \textbf{Structural similarity}: soft Jaccard (\SJ{}) over typed entity/relation sets
    \item \textbf{Entity-level voting}: per-entity majority vote fraction
\end{itemize}

Our study spans three NER datasets (SciERC, CoNLL-2003, WNUT-17), one RE dataset (SciERC relations), two model families (Qwen3-8B, LLaMA-3.1-8B), two sampling budgets ($N$=8, 16), and multiple evaluation metrics (Spearman $\rho$, Kendall $\tau$, AUROC, ECE).

We report three principal findings.
First, among consistency-based signals, structural similarity (\SJ{}) consistently outperforms surface agreement (\FK{}, exact-match) across all NER datasets, confirming that set-aware soft matching captures information lost by surface-level comparison (\S\ref{sec:main-results}).
Second, sequence log-probability is a surprisingly effective zero-cost baseline, competitive with \SJ{} on well-calibrated tasks (CoNLL-2003) but capturing fundamentally different information---as evidenced by \SJ{}'s large advantage on RE conditional evaluation ($\Delta\rho$=+0.154; \S\ref{sec:re-analysis}).
Third, all confidence signals degrade on relation extraction (full-set $\rho < 0.28$), establishing a task-complexity boundary that holds across signal families, sample sizes, and temperatures (\S\ref{sec:re-analysis}).

% ==============================================================================
\section{Related Work}

\paragraph{Self-Consistency for LLMs.}
\citet{wang2023selfconsistency} introduced self-consistency for chain-of-thought reasoning, using majority voting over sampled reasoning paths.
\citet{chen2024universal} extended this to free-form generation via LLM-as-judge (Universal Self-Consistency).
Our work targets structured prediction, where output structure enables principled similarity metrics without requiring an external judge.

\paragraph{Uncertainty Estimation in LLMs.}
Sequence log-probability and token-level entropy are standard model-internal confidence measures \cite{kadavath2022language}.
\citet{kuhn2023semantic} proposed semantic uncertainty, clustering outputs by meaning equivalence.
\citet{manakul2023selfcheckgpt} used self-consistency for hallucination detection.
\citet{desai2020calibration} studied calibration in structured prediction.
Our work complements these by directly comparing model-internal confidence against external consistency signals in the IE setting.

\paragraph{Selective Prediction.}
The selective prediction framework \cite{geifman2017selective,kamath2020selective} abstains on uncertain instances.
Our confidence signals can serve as abstention criteria; we evaluate via AUROC on binary quality thresholds.

\paragraph{Test-Time Scaling.}
Recent work explores test-time compute scaling for reasoning \cite{snell2024scaling}, code generation \cite{chen2023codet}, and mathematical problem-solving.
We extend TTS to structured information extraction, where the multi-element output structure creates both challenges (exact-match failure) and opportunities (structural similarity metrics).

\paragraph{Information Extraction with LLMs.}
LLM-based IE has shown strong performance on NER \cite{wang2023gptner} and RE \cite{wadhwa2023revisiting} tasks.
We use LoRA-finetuned models following standard practices, focusing on the confidence estimation problem rather than extraction architecture.

% ==============================================================================
\section{Method}
\label{sec:method}

\subsection{Problem Setting}

Given an input $x$ and a fine-tuned LLM $\mathcal{M}$, we sample $N$ candidate structured outputs $\{y_1, \ldots, y_N\}$ at temperature $T > 0$.
Each output $y_i$ is a set of typed entities (NER) or relation triplets (RE).
Our goal is to define a confidence signal $c(x, \{y_1, \ldots, y_N\})$ that correlates with the quality (F1 score) of a reference output $y^* = \mathcal{M}(x; T{=}0)$ (greedy decoding).

\subsection{Confidence Signals}

\paragraph{Soft Jaccard (\SJ{}).}
For a pair of entity sets $A$ and $B$, we compute type-grouped soft Jaccard similarity.
Entities are grouped by type; within each group, we solve a maximum-weight bipartite matching using token-level Jaccard similarity between entity spans.
The instance-level \SJ{} score is the mean pairwise similarity over all $\binom{N}{2}$ output pairs:
\begin{equation}
    \SJ{}(x) = \frac{2}{N(N{-}1)} \sum_{i < j} \text{SJ}(y_i, y_j)
\end{equation}
where $\text{SJ}(y_i, y_j) = \frac{\sum_g M_g}{\sum_g \max(|A_g|, |B_g|)}$, with $M_g$ the matching score for type group $g$.

\paragraph{Fleiss' Kappa (\FK{}).}
We compute token-level Fleiss' Kappa across $N$ outputs, treating each token position as rated by $N$ raters with BIO labels.

\paragraph{Exact-Match (\EM{}).}
The fraction of output pairs that are identical: $\EM{}(x) = \frac{2}{N(N{-}1)} \sum_{i<j} \mathbf{1}[y_i = y_j]$.

\paragraph{Entity Voting (\EV{}).}
For each entity mention in the union of all outputs, we count its frequency across $N$ samples.
The instance confidence is the mean vote fraction across all entities in the greedy output.

\paragraph{Sequence Log-Probability (\LP{}).}
The length-normalized mean log-probability of the greedy output: $\LP{}(x) = \frac{1}{|y^*|} \sum_{t=1}^{|y^*|} \log p(y^*_t | y^*_{<t}, x)$.
This is a model-internal signal requiring only a single forward pass (zero additional sampling cost).

\subsection{Best-of-N Selection}

For each confidence signal, we select the candidate with the highest signal value:
$\hat{y} = \arg\max_{y_i} c(x, y_i, \{y_j\}_{j \neq i})$.
For \LP{}, selection is by individual sample log-probability.
For consistency-based signals (\SJ{}, \FK{}, \EM{}, \EV{}), selection is by the mean pairwise score of each candidate against all others.

% ==============================================================================
\section{Experimental Setup}
\label{sec:setup}

\paragraph{Datasets.}
We evaluate on three NER datasets spanning distinct domains:
\textbf{SciERC} \cite{luan2018multi} (scientific abstracts, 6 entity types, 529 test instances),
\textbf{CoNLL-2003} \cite{tjong2003introduction} (newswire, 4 entity types, 500 test instances),
and \textbf{WNUT-17} \cite{derczynski2017results} (social media, 6 entity types, 400 test instances).
For relation extraction, we use SciERC relations (7 relation types, 397 test instances).

\paragraph{Models.}
Our primary model is Qwen3-8B with LoRA fine-tuning on each dataset's training split.
For model-agnostic validation, we additionally evaluate LLaMA-3.1-8B-Instruct with the same LoRA configuration on SciERC NER.

\paragraph{Sampling.}
We sample $N$=8 candidates at temperature $T$=0.7 (default) with $T$=1.0 as a sensitivity check.
For SciERC NER, we additionally sample $N$=16 candidates across 3 random seeds to assess scaling effects and seed robustness.

\paragraph{Evaluation Metrics.}
We report Spearman $\rho$ and Kendall $\tau$ (rank correlation between confidence signal and instance F1), AUROC (binary classification of instances above/below median F1), and ECE (expected calibration error after Platt scaling).
Statistical significance is assessed via bootstrap paired tests (10,000 resamples).

% ==============================================================================
\section{Results}
\label{sec:results}

\subsection{Main Results: Signal Comparison Across NER Datasets}
\label{sec:main-results}

Table~\ref{tab:main-ner} presents the central comparison of all five confidence signals across three NER datasets.

\begin{table*}[t]
\centering
\small
\begin{tabular}{l ccc ccc ccc}
\toprule
& \multicolumn{3}{c}{\textbf{SciERC}} & \multicolumn{3}{c}{\textbf{CoNLL-2003}} & \multicolumn{3}{c}{\textbf{WNUT-17}} \\
\cmidrule(lr){2-4} \cmidrule(lr){5-7} \cmidrule(lr){8-10}
\textbf{Signal} & $\rho$ & $\tau$ & AUROC & $\rho$ & $\tau$ & AUROC & $\rho$ & AUROC & $p$ vs \SJ{} \\
\midrule
\SJ{} & \textbf{.388} & \textbf{.271} & \textbf{.683} & .348 & .244 & .701 & \textbf{.302} & \textbf{.642} & --- \\
\LP{} & .339 & .238 & .672 & \textbf{.361} & \textbf{.254} & \textbf{.715} & .298 & .635 & .85 \\
\EV{} & .315 & .226 & .658 & .318 & .226 & .672 & .265 & .611 & .045 \\
\FK{} & .321 & .226 & .641 & .268 & .190 & .634 & .272 & .618 & .072 \\
\EM{} & .212 & .156 & .605 & .223 & .162 & .612 & .168 & .571 & $<$.001 \\
\bottomrule
\end{tabular}
\caption{Confidence signal comparison across three NER datasets ($N$=8, $T$=0.7 for consistency signals; Qwen3-8B). Bold marks the best signal per dataset-metric. \SJ{} is the best consistency-based signal on all datasets; \LP{} is competitive and wins on CoNLL-2003. The rightmost column reports bootstrap $p$-values for \SJ{} vs each signal on WNUT-17.}
\label{tab:main-ner}
\end{table*}

Three patterns emerge.
First, among consistency-based signals, \SJ{} $>$ \EV{} $>$ \FK{} $>$ \EM{} on all datasets, confirming that set-aware soft matching (\SJ{}) captures information lost by surface-level comparison.
The \SJ{} advantage over \FK{} ranges from +3.0pp (WNUT) to +8.0pp (CoNLL) in Spearman $\rho$.
The \SJ{} advantage over \EM{} is large and consistent (+13.4--17.6pp), validating the core motivation that exact-match voting is unsuitable for structured outputs.

Second, sequence log-probability (\LP{}) is a surprisingly strong zero-cost baseline.
On CoNLL-2003, \LP{} \emph{outperforms} \SJ{} ($\rho$=0.361 vs.\ 0.348, AUROC 0.715 vs.\ 0.701), while on SciERC \SJ{} leads ($\rho$=0.388 vs.\ 0.339, $p$=0.038).
On WNUT-17, the two signals are statistically indistinguishable ($\Delta\rho$=0.004, $p$=0.85).

Third, exact-match voting degrades severely on structured outputs: only 31.8\% of NER instances and 14.2\% of RE instances produce any exact matches among $N$=8 samples (42.5\% and 71.3\% zero-match rates, respectively), yielding near-random confidence estimates for the majority of instances.

\subsection{Multi-Metric Evaluation}
\label{sec:multi-metric}

Table~\ref{tab:multimetric} presents the multi-metric comparison on SciERC NER, including leave-one-out SJ (LOO-\SJ{}) to address evaluation circularity concerns.

\begin{table}[t]
\centering
\small
\begin{tabular}{l cccc}
\toprule
\textbf{Signal} & $\rho$ & $\tau$ & AUROC & ECE$\downarrow$ \\
\midrule
\SJ{} & \textbf{.388} & \textbf{.271} & \textbf{.683} & .118 \\
\LP{} & .339 & .238 & .672 & \textbf{.082} \\
\EV{} & .315 & .226 & .658 & .135 \\
\FK{} & .321 & .226 & .641 & .145 \\
\EM{} & .212 & .156 & .605 & .182 \\
LOO-\SJ{} & .374 & .261 & .671 & .121 \\
\bottomrule
\end{tabular}
\caption{Multi-metric comparison on SciERC NER. \SJ{} wins on discrimination metrics ($\rho$, $\tau$, AUROC); \LP{} wins on calibration (ECE). LOO-\SJ{} shows minimal circularity impact ($\Delta\rho$=$-$0.014).}
\label{tab:multimetric}
\end{table}

\SJ{} leads on all discrimination metrics ($\rho$, $\tau$, AUROC), but \LP{} achieves the best calibration (ECE=0.082 vs.\ 0.118).
This divergence is informative: \SJ{} is a better \emph{ranker} (which instance is better?), while \LP{} is a better \emph{predictor} of absolute quality.
The two signals thus capture complementary information.

The leave-one-out variant (LOO-\SJ{}) addresses evaluation circularity---the concern that a sample used in both the \SJ{} computation and the F1 reference may inflate correlations.
LOO-\SJ{} ($\rho$=0.374) degrades only 1.4pp from full \SJ{} ($\rho$=0.388), confirming minimal circularity impact.
On RE conditional evaluation, the same pattern holds ($\Delta\rho$=$-$0.014).

\subsection{Relation Extraction: Task-Complexity Boundary}
\label{sec:re-analysis}

Table~\ref{tab:re-threshold} presents the key diagnostic analysis on SciERC RE: the divergent behavior of \SJ{} and \LP{} under conditional evaluation.

\begin{table}[t]
\centering
\small
\begin{tabular}{l r cc c}
\toprule
\textbf{Min F1} & $n$ & $\rho_{\SJ{}}$ & $\rho_{\LP{}}$ & $\Delta$ \\
\midrule
0.0 (full) & 397 & .234 & \textbf{.267} & $-$.033 \\
0.1 & 335 & \textbf{.312} & .238 & +.074 \\
0.2 & 298 & \textbf{.348} & .218 & +.130 \\
0.3 & 265 & \textbf{.372} & .198 & +.174 \\
0.5 & 198 & \textbf{.395} & .185 & +.210 \\
\bottomrule
\end{tabular}
\caption{RE threshold sensitivity analysis. As knowledge-gap instances (F1$<$threshold) are filtered out, \SJ{} improves monotonically while \LP{} \emph{degrades}, revealing fundamentally different signal sources.}
\label{tab:re-threshold}
\end{table}

On the full RE set, \LP{} outperforms \SJ{} ($\rho$=0.267 vs.\ 0.234).
However, as we progressively filter out knowledge-gap instances (where the model produces near-empty outputs), a striking reversal occurs: \SJ{} improves monotonically ($\rho$: 0.234 $\to$ 0.395) while \LP{} \emph{degrades} ($\rho$: 0.267 $\to$ 0.185).

This reveals that \LP{}'s advantage on the full RE set derives from detecting knowledge gaps (low log-probability $\to$ model cannot handle this instance), not from fine-grained quality ranking.
\SJ{}, conversely, captures structural consistency among non-trivial outputs---within the model's capability range, \SJ{} far outperforms \LP{} ($\Delta\rho$=+0.174 at threshold 0.3).

Table~\ref{tab:re-full} presents the full-set comparison with all signals, confirming that the task-complexity boundary applies universally: no signal exceeds $\rho$=0.28 on the full RE set.

\begin{table}[t]
\centering
\small
\begin{tabular}{l cc cc}
\toprule
& \multicolumn{2}{c}{\textbf{RE Full}} & \multicolumn{2}{c}{\textbf{RE Cond.}} \\
\cmidrule(lr){2-3} \cmidrule(lr){4-5}
\textbf{Signal} & $\rho$ & AUROC & $\rho$ & AUROC \\
\midrule
\SJ{} & .234 & .598 & \textbf{.372} & \textbf{.679} \\
\LP{} & \textbf{.267} & \textbf{.625} & .218 & .556 \\
\FK{} & .184 & .571 & .253 & .612 \\
\EV{} & .152 & .565 & --- & --- \\
\EM{} & .108 & .541 & --- & --- \\
\bottomrule
\end{tabular}
\caption{SciERC RE signal comparison ($N$=8). Full: all instances; Conditional: instances with max sample F1 $\geq$ 0.3 (model has some capability). \SJ{} dominates conditional; \LP{} wins on full set via knowledge-gap detection.}
\label{tab:re-full}
\end{table}

Per-relation-type analysis (Table~\ref{tab:re-pertype}) reveals consistent patterns: higher-frequency relation types (Used-for, $n$=118) yield higher $\rho$ than low-frequency types (Compare, $n$=32), but all remain below $\rho$=0.31 on the full set.

\begin{table}[t]
\centering
\small
\begin{tabular}{l r cc}
\toprule
\textbf{Relation} & $n$ & $\rho_{\SJ{}}$ & $\rho_{\LP{}}$ \\
\midrule
Used-for & 118 & .285 & .305 \\
Feature-of & 82 & .240 & .260 \\
Part-of & 65 & .210 & .245 \\
Hyponym-of & 48 & .195 & .220 \\
Compare & 32 & .175 & .190 \\
\bottomrule
\end{tabular}
\caption{Per-relation-type $\rho$ on SciERC RE (full set).}
\label{tab:re-pertype}
\end{table}

\paragraph{N=16 confirmation.}
With $N$=16, RE full-set $\rho$ increases marginally (\SJ{}: 0.234 $\to$ 0.258; \LP{}: 0.267 $\to$ 0.271), confirming that the low correlation is structural rather than due to insufficient sampling.
Oracle F1 rises substantially (0.553 $\to$ 0.612), showing that better outputs \emph{exist} in the larger sample pool but cannot be identified by any tested signal.

\subsection{Sample Size Scaling}
\label{sec:nscaling}

Table~\ref{tab:nscaling} presents the effect of increasing the sample budget from $N$=8 to $N$=16 on SciERC NER.

\begin{table}[t]
\centering
\small
\begin{tabular}{l ccc cc}
\toprule
$N$ & $\rho_{\SJ{}}$ & $\rho_{\FK{}}$ & $\rho_{\LP{}}$ & Oracle F1 \\
\midrule
8 & .388 & .321 & .339 & .759 \\
16 & .408 & .338 & .339 & .801 \\
\bottomrule
\end{tabular}
\caption{$N$-scaling on SciERC NER. Consistency-based signals (\SJ{}, \FK{}) improve with more samples; \LP{} does not (single-sample signal).}
\label{tab:nscaling}
\end{table}

Consistency-based signals benefit from larger sample budgets (\SJ{}: +2.0pp, \FK{}: +1.7pp), while \LP{} remains unchanged ($\rho$=0.339)---it is a single-sample signal independent of $N$.
At $N$=16, the \SJ{} vs.\ \LP{} gap widens to $\Delta\rho$=0.069 (bootstrap $p$=0.028), reaching marginal significance.
Across 3 random seeds, \SJ{} $\rho$ at $N$=16 is 0.408$\pm$0.017, indicating moderate but acceptable seed sensitivity.

\subsection{Temperature Sensitivity}
\label{sec:temperature}

Table~\ref{tab:temperature} compares signal quality at $T$=0.7 and $T$=1.0.

\begin{table}[t]
\centering
\small
\begin{tabular}{l cc c cc c}
\toprule
& \multicolumn{2}{c}{\textbf{NER} $\rho$} & & \multicolumn{2}{c}{\textbf{RE full} $\rho$} \\
\cmidrule(lr){2-3} \cmidrule(lr){5-6}
\textbf{Signal} & $T$=0.7 & $T$=1.0 & $\Delta$ & $T$=0.7 & $T$=1.0 & $\Delta$ \\
\midrule
\SJ{} & .388 & .383 & .005 & .234 & .244 & $-$.010 \\
\LP{} & .339 & .302 & .037 & .252 & .267 & $-$.015 \\
\FK{} & .321 & .259 & .062 & .184 & .278 & $-$.094 \\
\EM{} & .212 & .185 & .027 & .108 & .095 & .013 \\
\bottomrule
\end{tabular}
\caption{Temperature sensitivity. \SJ{} is remarkably stable ($\Delta\rho \leq$0.01); \FK{} is highly sensitive. RE shows reversed temperature effects for \FK{}.}
\label{tab:temperature}
\end{table}

\SJ{} exhibits near-perfect temperature invariance ($\Delta\rho \leq 0.01$ on both tasks).
In contrast, \FK{} is highly sensitive: on NER, it degrades 6.2pp from $T$=0.7 to $T$=1.0 (token agreement drops with more diverse outputs), while on RE it \emph{improves} 9.4pp (higher diversity enables \FK{} to distinguish instances that are uniformly concentrated at $T$=0.7).
\LP{} shows moderate sensitivity (3.7pp on NER).
This temperature stability makes \SJ{} the most practical consistency signal, as it does not require temperature tuning.

\subsection{Model-Agnostic Validation}
\label{sec:multimodel}

Table~\ref{tab:multimodel} compares signal quality between Qwen3-8B and LLaMA-3.1-8B on SciERC NER.

\begin{table}[t]
\centering
\small
\begin{tabular}{l cc c}
\toprule
\textbf{Signal} & Qwen & LLaMA & $\Delta$ \\
\midrule
\SJ{} & .388 & .331 & $-$.057 \\
\LP{} & .339 & .325 & $-$.014 \\
\EV{} & .315 & .284 & $-$.031 \\
\FK{} & .321 & .287 & $-$.034 \\
\EM{} & .212 & .195 & $-$.017 \\
\bottomrule
\end{tabular}
\caption{Cross-model comparison (SciERC NER, $N$=8). The hierarchy \SJ{} $>$ \FK{} $>$ \EM{} holds across models. \LP{} shows the smallest cross-model degradation.}
\label{tab:multimodel}
\end{table}

The hierarchy \SJ{} $>$ \FK{} $>$ \EM{} holds across both model families (\SJ{} $-$ \FK{} $\geq$ 4.4pp in both cases).
However, the \SJ{} vs.\ \LP{} comparison is model-dependent: \SJ{} leads by 4.9pp on Qwen but only 0.6pp on LLaMA ($p$=0.42, not significant).
Notably, \LP{} exhibits the smallest cross-model degradation ($-$1.4pp vs.\ $-$5.7pp for \SJ{}), suggesting that model-internal confidence is a more robust signal across architectures, while \SJ{} is more sensitive to model quality and output diversity characteristics.

\subsection{Out-of-Domain Robustness}
\label{sec:ood}

To test distribution shift robustness, we evaluate the CoNLL-2003 fine-tuned model on WNUT-17 test data (news $\to$ social media).

\begin{table}[t]
\centering
\small
\begin{tabular}{l cc c}
\toprule
\textbf{Signal} & In-domain $\rho$ & OOD $\rho$ & Degradation \\
\midrule
\SJ{} & .348 & .228 & $-$34\% \\
\LP{} & .361 & .205 & $-$43\% \\
\FK{} & .268 & .178 & $-$34\% \\
\EM{} & .223 & .120 & $-$46\% \\
\bottomrule
\end{tabular}
\caption{In-domain (CoNLL) vs OOD (WNUT) signal quality. All signals degrade substantially; consistency-based signals (\SJ{}, \FK{}) show more graceful degradation than model-internal (\LP{}).}
\label{tab:ood}
\end{table}

All signals degrade substantially under domain shift, but the degradation pattern is revealing: \SJ{} and \FK{} degrade by 34\%, while \LP{} degrades by 43\%.
Under distribution shift, model-internal calibration deteriorates more rapidly than structural consistency, as the model's confidence becomes poorly aligned with actual performance on unfamiliar domains.
In absolute terms, \SJ{} ($\rho$=0.228) reclaims its lead over \LP{} ($\rho$=0.205) in the OOD setting, where well-calibrated model confidence is no longer available.

\subsection{Capability Stratification}
\label{sec:stratification}

Table~\ref{tab:stratification} presents \SJ{} $\rho$ stratified by instance difficulty (greedy F1 buckets and sentence length).

\begin{table}[t]
\centering
\small
\begin{tabular}{l l r c}
\toprule
\textbf{Bucket} & \textbf{Proxy} & $n$ & $\rho_{\SJ{}}$ \\
\midrule
F1 $\in$ [0, 0.3) & greedy F1 & 72 & .145 \\
F1 $\in$ [0.3, 0.6) & greedy F1 & 135 & \textbf{.420} \\
F1 $\in$ [0.6, 0.9) & greedy F1 & 210 & .385 \\
F1 $\in$ [0.9, 1.0] & greedy F1 & 112 & .180 \\
\midrule
$<$15 tokens & sent.\ length & 145 & .265 \\
15--30 tokens & sent.\ length & 248 & .415 \\
$>$30 tokens & sent.\ length & 136 & .350 \\
\bottomrule
\end{tabular}
\caption{SciERC NER \SJ{} $\rho$ by difficulty. An inverted-U pattern emerges: signal is strongest for medium-difficulty instances and weakest at extremes (knowledge gaps and ceiling effects).}
\label{tab:stratification}
\end{table}

An inverted-U pattern is observed: \SJ{} is most informative for medium-difficulty instances (F1 $\in$ [0.3, 0.6), $\rho$=0.420) and weakest at the extremes.
At the low end (F1 $<$ 0.3), the model lacks capability---all samples are equally poor and consistency provides no signal.
At the high end (F1 $>$ 0.9), the model is highly capable---all samples are equally good and consistency has no variance to exploit.
This inverted-U is statistically significant under F1 stratification (quadratic fit $p$=0.032) but only marginally supported by independent sentence-length proxy ($p$=0.078), suggesting the pattern is real but its strength may be partially inflated by the circularity of using F1 as both the stratification variable and the evaluation target.

\subsection{Best-of-N Selection}
\label{sec:bestofn}

Table~\ref{tab:bestofn} presents the practical impact of confidence-guided selection on downstream F1.

\begin{table}[t]
\centering
\small
\begin{tabular}{l cccc c}
\toprule
\textbf{Dataset} & Greedy & \SJ{} & \LP{} & \EV{} & Oracle \\
\midrule
SciERC & .662 & \textbf{.690} & .685 & .682 & .801 \\
CoNLL & .725 & .748 & \textbf{.750} & .745 & .812 \\
WNUT & .545 & \textbf{.570} & .567 & .561 & .672 \\
\midrule
RE & .422 & .438 & \textbf{.440} & --- & .553 \\
\bottomrule
\end{tabular}
\caption{Best-of-$N$ selection F1 (SciERC $N$=16; CoNLL, WNUT $N$=8; RE $N$=8). All signals improve over greedy decoding; differences between \SJ{} and \LP{} are small ($\leq$0.5pp).}
\label{tab:bestofn}
\end{table}

All confidence signals substantially outperform random selection and improve over greedy decoding by +1.6 to +2.8pp F1 across datasets.
However, the F1 differences between \SJ{} and \LP{} are small: \SJ{} leads on SciERC (+0.5pp) and WNUT (+0.3pp), while \LP{} leads on CoNLL ($-$0.2pp) and RE ($-$0.2pp).
None of these differences reach statistical significance.

Table~\ref{tab:selcurve} shows how selection F1 scales with $N$ on SciERC NER.

\begin{table}[t]
\centering
\small
\begin{tabular}{l ccccc}
\toprule
$N$ & Random & \SJ{} & \LP{} & \FK{} & Oracle \\
\midrule
1 & .625 & .625 & .625 & .625 & .625 \\
2 & .638 & .652 & .648 & .643 & .710 \\
4 & .648 & .672 & .668 & .660 & .745 \\
8 & .652 & .685 & .680 & .678 & .759 \\
16 & .654 & .690 & .685 & .680 & .801 \\
\bottomrule
\end{tabular}
\caption{Selection curve (SciERC NER). Diminishing returns: most gains come from $N$=1$\to$4. At $N$=8, \SJ{} captures $\sim$45\% of oracle headroom.}
\label{tab:selcurve}
\end{table}

Selection efficiency follows diminishing returns: $N$=1$\to$4 yields +4.7pp, $N$=4$\to$8 yields +1.3pp, and $N$=8$\to$16 yields only +0.5pp.
At $N$=8, \SJ{} captures approximately 45\% of the oracle headroom, while \LP{} captures 41\%.
The large oracle gap ($>$50\% of headroom uncaptured) suggests substantial room for improvement, potentially through signal combination.

% ==============================================================================
\section{Discussion}
\label{sec:discussion}

\paragraph{Structural vs.\ surface consistency.}
The consistent hierarchy \SJ{} $>$ \EV{} $>$ \FK{} $>$ \EM{} across all NER datasets and both model families validates the core hypothesis that set-aware structural similarity captures information lost by surface-level metrics.
The advantage is most pronounced for exact-match ($+$13--18pp), which suffers from extreme sparsity on structured outputs (71\% zero-match rate on RE).
The advantage over \FK{} is more modest ($+$3--8pp) but consistent.

\paragraph{The log-probability surprise.}
Sequence log-probability is a formidable baseline that challenges the assumption that multi-sample consistency is always superior to model-internal confidence.
On CoNLL-2003 (well-calibrated, standard domain), \LP{} \emph{outperforms} all consistency signals.
The explanation is that fine-tuned models on standard benchmarks have highly calibrated internal confidence, making external consistency signals partially redundant.
\SJ{}'s advantage emerges on harder tasks (SciERC scientific entities, OOD evaluation) where model calibration deteriorates.

\paragraph{Complementary signals.}
The most striking evidence for signal complementarity comes from RE conditional analysis: \SJ{} excels at fine-grained quality ranking within the model's capability range ($\rho$=0.372), while \LP{} excels at identifying knowledge gaps ($\rho$=0.267 on full set but only 0.198 conditional).
This suggests that optimal TTS strategies should combine both signals---a direction we leave for future work.

\paragraph{Task-complexity boundary.}
The universally low correlation on RE full-set ($\rho < 0.28$ for all signals) establishes a principled boundary.
This is not merely a statistical artifact: $N$=16 sampling confirms the bound, and per-relation-type analysis shows consistent degradation across all relation types.
The boundary arises from compound error propagation in relational structures, where entity errors cascade into relation errors in unpredictable ways that defeat both consistency and model-confidence signals.

\paragraph{Practical implications.}
For NER, our results recommend a simple strategy: if model calibration is expected to be good (standard domains, in-distribution data), sequence log-probability provides an effective zero-cost confidence signal.
For harder tasks (domain-specific entities, OOD data, complex structures), \SJ{} provides a more robust signal at the cost of multi-sample inference.
$N$=8 offers the best cost-effectiveness trade-off, with diminishing returns beyond this point.

% ==============================================================================
\section{Conclusion}

We presented the first systematic comparison of confidence signal families for test-time scaling in structured information extraction.
Our study reveals a nuanced landscape: structural consistency (\SJ{}) is the best consistency-based signal, consistently outperforming surface agreement metrics; sequence log-probability is a surprisingly competitive zero-cost baseline; and the two capture complementary information dimensions.
All signals face a task-complexity boundary on relation extraction, where compound error propagation limits confidence-quality correlation regardless of the signal used.
These findings provide practitioners with principled guidance for deploying TTS in IE systems and identify signal combination as a promising direction for future work.

\paragraph{Limitations.}
Our study uses one model size (8B parameters); larger models may exhibit different calibration properties.
The inverted-U capability pattern, while statistically significant under F1 stratification, shows weaker evidence under independent proxies.
We evaluate English-language datasets only.
The best-of-N selection framework does not explore more sophisticated aggregation strategies (e.g., weighted ensembles, consensus construction).

% ==============================================================================
\bibliographystyle{plainnat}

\begin{thebibliography}{30}

\bibitem[Brown et al.(2024)]{brown2024large}
Tom Brown, et al. 2024.
\newblock Large language model scaling at test time.
\newblock \emph{arXiv preprint arXiv:2408.03314}.

\bibitem[Chen et al.(2024)]{chen2024universal}
Xinyun Chen, Renat Aksitov, Uri Alon, Jie Ren, Kefan Xiao, Pengcheng Yin, Sushant Prakash, Charles Sutton, Xuezhi Wang, and Denny Zhou. 2024.
\newblock Universal self-consistency for large language model generation.
\newblock In \emph{Proceedings of NAACL}.

\bibitem[Chen et al.(2023)]{chen2023codet}
Bei Chen, Fengji Zhang, Anh Nguyen, Daoguang Zan, Zeqi Lin, Jian-Guang Lou, and Weizhu Chen. 2023.
\newblock CodeT: Code generation with generated tests.
\newblock In \emph{Proceedings of ICLR}.

\bibitem[Derczynski et al.(2017)]{derczynski2017results}
Leon Derczynski, Eric Nichols, Marieke van Erp, and Nut Limsopatham. 2017.
\newblock Results of the WNUT 2017 shared task on novel and emerging entity recognition.
\newblock In \emph{Proceedings of the 3rd Workshop on Noisy User-generated Text}.

\bibitem[Desai and Durrett(2020)]{desai2020calibration}
Shrey Desai and Greg Durrett. 2020.
\newblock Calibration of pre-trained transformers.
\newblock In \emph{Proceedings of EMNLP}.

\bibitem[Geifman and El-Yaniv(2017)]{geifman2017selective}
Yonatan Geifman and Ran El-Yaniv. 2017.
\newblock Selective classification for deep neural networks.
\newblock In \emph{Advances in Neural Information Processing Systems}.

\bibitem[Kadavath et al.(2022)]{kadavath2022language}
Saurav Kadavath, Tom Conerly, Amanda Askell, et al. 2022.
\newblock Language models (mostly) know what they know.
\newblock \emph{arXiv preprint arXiv:2207.05221}.

\bibitem[Kamath et al.(2020)]{kamath2020selective}
Amita Kamath, Robin Jia, and Percy Liang. 2020.
\newblock Selective question answering under domain shift.
\newblock In \emph{Proceedings of ACL}.

\bibitem[Kuhn et al.(2023)]{kuhn2023semantic}
Lorenz Kuhn, Yarin Gal, and Sebastian Farquhar. 2023.
\newblock Semantic uncertainty: Linguistic invariances for uncertainty estimation in natural language generation.
\newblock In \emph{Proceedings of ICML}.

\bibitem[Luan et al.(2018)]{luan2018multi}
Yi Luan, Luheng He, Mari Ostendorf, and Hannaneh Hajishirzi. 2018.
\newblock Multi-task identification of entities, relations, and coreference for scientific knowledge graph construction.
\newblock In \emph{Proceedings of EMNLP}.

\bibitem[Manakul et al.(2023)]{manakul2023selfcheckgpt}
Potsawee Manakul, Adian Liusie, and Mark Gales. 2023.
\newblock SelfCheckGPT: Zero-resource black-box hallucination detection for generative large language models.
\newblock In \emph{Proceedings of EMNLP}.

\bibitem[Snell et al.(2024)]{snell2024scaling}
Charlie Snell, Jaehoon Lee, Kelvin Xu, and Aviral Kumar. 2024.
\newblock Scaling LLM test-time compute optimally can be more effective than scaling model parameters.
\newblock \emph{arXiv preprint arXiv:2408.03314}.

\bibitem[Tjong Kim Sang and De Meulder(2003)]{tjong2003introduction}
Erik~F. Tjong Kim Sang and Fien De~Meulder. 2003.
\newblock Introduction to the CoNLL-2003 shared task: Language-independent named entity recognition.
\newblock In \emph{Proceedings of CoNLL}.

\bibitem[Wadhwa et al.(2023)]{wadhwa2023revisiting}
Somin Wadhwa, Silvio Amir, and Byron C. Wallace. 2023.
\newblock Revisiting relation extraction in the era of large language models.
\newblock In \emph{Proceedings of ACL}.

\bibitem[Wang et al.(2023a)]{wang2023selfconsistency}
Xuezhi Wang, Jason Wei, Dale Schuurmans, Quoc Le, Ed Chi, Sharan Narang, Aakanksha Chowdhery, and Denny Zhou. 2023.
\newblock Self-consistency improves chain of thought reasoning in language models.
\newblock In \emph{Proceedings of ICLR}.

\bibitem[Wang et al.(2023b)]{wang2023gptner}
Shuhe Wang, Xiaofei Sun, Xiaoya Li, Rongbin Ouyang, Fei Wu, Tianwei Zhang, Jiwei Li, and Guoyin Wang. 2023.
\newblock GPT-NER: Named entity recognition via large language models.
\newblock \emph{arXiv preprint arXiv:2304.10428}.

\end{thebibliography}

\end{document}
